Robot manipulation is often framed as a geometric reasoning problem: identify a grasp, plan a collision-free path through configuration space, and move the object to its goal. 
This abstraction has been enormously productive because it converts manipulation into problems that are comparatively well understood---reachability, collision avoidance, grasp stability, and trajectory tracking. 
Yet it also hides much of what makes physical interaction powerful.
However, many useful manipulation behaviors are defined
% Many useful manipulation behaviors are not defined by static configurations alone, but
by dynamically feasible transitions between them: an object displaced by an impulsive contact, a liquid surface stabilized during transport, a damaged robot arm repurposed for whole-body contact, or a heavy payload moved near the limits of actuation.

This thesis has argued that these behaviors require treating dynamics not as a residual artifacts to be corrected, but as structured information to be represented, reused, and exploited during motion generation. 
Across the chapters that follow from this argument, dynamic feasibility is represented at increasingly amortized levels: first through online kinodynamic search with physics simulation, then through reusable motion structure for multi-query planning, and finally through learned conditional models that generate dynamically informed trajectories and optimization seeds. 
This chapter summarizes the main contributions, revisits them through the four design axes introduced in \cref{ch:framework}, and identifies open directions for building manipulation systems that operate closer to the physical limits of their embodiments.

% Robot manipulation research has long defaulted to a geometric view of the problem: plan a collision-free path through configuration space, close the gripper, and move the object to its goal.
% This is computationally convenient---once an object is under grasp, its state becomes predictable and control straightforward.
% But grasping is inherently limited by robot reachability, end-effector design, and the physical properties of the objects themselves.
% More broadly, the full range of tasks that manipulation actually requires---poking objects into position, transporting liquid-filled containers without spilling, operating with a damaged joint, handling payloads beyond rated capacity---depends critically on understanding and exploiting the dynamics of the task, not just its geometry.

% This thesis has argued that explicitly incorporating dynamic constraints into the motion generation process allows robots to perform a substantially broader range of manipulation tasks while operating closer to their true physical limits.
% The work spans three thematic areas: \textit{sampling-based kinodynamic planning}, which integrates dynamics directly into the search process through physics simulation; \textit{generative models for dynamics-aware motion}, which encode dynamic constraints implicitly through conditioning; and \textit{dynamics-informed trajectory seeding}, which improves the convergence of constrained optimization by initializing near the constraint manifold.
% This chapter summarizes what each contribution established, reflects on the collective arc through the lens of the four design axes introduced in \cref{ch:framework}, and identifies the most promising directions for future work.